% Agent 27: Pages 447--450 (PDF pp.52--54)
% Sections 6.2 The Growth of Cosmological Holes by Accretion
%          6.3 Limit on the Number of Baryons in Cosmological Holes
%          6.4 Caution
%          References (beginning)
% =============================================================================

%% ----------------------------------------------------------------
%% 6.2  The Growth of Cosmological Holes by Accretion
%% ----------------------------------------------------------------

\subsection{The Growth of Cosmological Holes by Accretion}
\label{sec:6.2}

If cosmological holes have existed since near the beginning of the universe,
then during the era when the energy density in radiation exceeded that in matter,
accretion of radiation onto the holes must have been important (Zel'dovich and
Novikov 1966).  For a rough estimate of the effects of such accretion we can use
equation~(4.2.14) for the accretion of gas onto a stationary hole, taking the
velocity of sound to be $a_\infty = c/\sqrt{3}$:
%
\begin{equation}
\tag{6.2.1}
dM/dt \simeq r_g^2 \rho^{\text{rad}}\!.
\end{equation}
%
In standard cosmological models the radiation energy density varies as
%
\begin{equation}
\tag{6.2.2}
\rho^{\text{rad}} = 1/Gt^2.
\end{equation}

Putting this into the growth equation~(6.2.1), and integrating from the moment
$t_0$ when a hole of initial mass $M_0$ first forms, we obtain for the final mass, at
$t > t_0$:
%
\begin{equation}
\tag{6.2.3}
M = M_0 (1 - G M_0 / c^3 t_0)^{-1}.
\end{equation}

Notice that the final mass diverges for $t_0 \to GM_0/c^3$.  However, the method of
calculation breaks down for this same limit because, for $t_0 \sim GM_0/c^3$ the
characteristic time scale for the accretion is the same as that for the change of
$\rho^{\text{rad}}$, so the ``steady-state'' hypothesis on which equation
(6.2.1) is based.  To determine whether the accretion is catastrophically great at
early times ($t_0 \lesssim GM_0/c^3$), one must solve the accretion problem in a nonsteady,
cosmological context.  One would not be surprised if the final answer depends
sensitively on the particular choice of initial conditions ($t_0 = GM_0/c^3$ or
$t_0 = 0.01\;GM_0/c^3$ or $t_0 = 10^{-10}\,GM_0/c^3$).  But the problem has not yet been
solved.

%% ----------------------------------------------------------------
%% 6.3  Limit on the Number of Baryons in Cosmological Holes
%% ----------------------------------------------------------------

\subsection{Limit on the Number of Baryons in Cosmological Holes}
\label{sec:6.3}

Independently of the issue of accretion, one can place a tight limit on what
fraction~$\alpha$ of all baryons in the Universe were swallowed into cosmological holes
in the early stages.

Suppose that prior to some particular early moment~$t_*$, a fraction~$\alpha$ of all
baryons had been swallowed into holes.  At the moment~$t_*$, the ratio of rest mass-energy in baryons to mass-energy in radiation and pairs was tiny
%
\begin{equation}
\tag{6.3.1}
\beta \equiv \left[\frac{\rho_0}{\rho^{\text{rad}} + \rho^{\text{(pairs)}}}\right]_{t_*}
\sim 10^{-7} (t_*/1\;\text{sec})^{1/2} \ll 1.
\end{equation}
%
However, as the universe subsequently expanded, this ratio increased until today
it is far greater than unity.  Moreover, the total mass-energy in holes divided by
the volume of the universe ($\alpha \rho_0$) changed in the same manner as
$\rho_0$ ($\propto$~vol$^{-1}$), if no new holes were formed after~$t_*$, and if accretion was negligible.
Otherwise it increased relative to~$\rho_0$.  This means that
%
\begin{equation}
\tag{6.3.2}
\frac{\alpha}{\beta(1-\alpha)}
\leq \left[\frac{\rho^{\text{(holes)}}}{\rho_0}\right]_{t_*}
\leq \left[\frac{\rho^{\text{(holes)}}}{\rho_0}\right]_{\text{today}}
< 80
\end{equation}
%
Here $\alpha/[\beta(1-\alpha)]$ is the value of $[\rho^{\text{(holes)}}/\rho_0]_{t_*}$ if the holes were all created from
primeval plasma at the moment~$t_*$; if they were created even earlier, when rest
mass was even less important, $[\rho^{\text{(holes)}}/\rho_0]_*$ would be even bigger.  The number
80 is a generous observational upper limit on the ratio of nonluminous matter
to luminous matter in the universe today.

Equation~(6.3.2) for $\beta \ll 1/80$ can be rewritten as
%
\begin{equation}
\tag{6.3.3}
\alpha < 1/80\beta \sim 10^{-5} (t_*/1\;\text{sec})^{-1/2}\,.
\end{equation}

This is a very tight limit on the amount of rest mass that could have been down
black holes at early times~$t_*$.

One can also place a tight limit on the amount of mass-energy that white holes
have spewed forth as radiation in recent times (at cosmological redshifts $z \lesssim 10$).
That mass-energy cannot exceed the total mass-energy in radiation today,
%
\begin{equation}
\tag{6.3.4}
\rho_{\text{from white holes}} < \rho_0^{\text{(rad)}}
= 5 \times 10^{-13}\;\text{erg/cm}^3
= 6 \times 10^{-34}\;\text{g/cm}^3\,.
\end{equation}

%% ----------------------------------------------------------------
%% 6.4  Caution
%% ----------------------------------------------------------------

\subsection{Caution}
\label{sec:6.4}

We conclude with a word of caution.  This cosmological section has dealt with
objects about which both theory and observation are equivocal.  There is no firm
theoretical reason for believing that cosmological holes exist, though theory
surely permits them.

By contrast, theory virtually demands the existence of black holes formed
by stellar collapse.  It would be astonishing indeed if every star in the Galaxy
somehow managed to avoid the black-hole fate!  [See Hoyle, Fowler, Burbidge,
and Burbidge (1964).]

The clear discovery of black holes, we expect, is only a few months or years
away.  (The X-ray source Cyg~X-1 is an excellent candidate.)  But cosmological
holes might never be found and might, in fact, not exist at all.  On the other
hand, we might eventually discover that they are the energizers of quasars and
of explosions in galactic nuclei.

%% ----------------------------------------------------------------
%% References
%% ----------------------------------------------------------------

\begin{thebibliography}{99}

\bibitem{Allen1963}
Allen, L.~H. 1963, \textit{Astrophysics---The Atmosphere of the Sun and Stars} (New York:
The Ronald Press Company).

\bibitem{Ambartsumyan1960}
Ambartsumyan, V.~A. 1960, \textit{Collected Works}, 2 (Yerevan).

\bibitem{Ambartsumyan1961}
Ambartsumyan, V.~A. 1961, 1964, \textit{Rapport 13 Conseil\`e Physique Solvay} (Brussels: Stoops).

\bibitem{Arnett1969}
Arnett, W.~D. 1969, \textit{Ap.\ Space Sci.}\ \textbf{5}, 180.

\bibitem{Arnett1973}
Arnett, W.~D. 1973, \textit{Ap.\ J.\ Space Sci.}\ \textbf{5}, 180.

\bibitem{Bardeen1970}
Bardeen, J.~M. 1970, \textit{Ap.\ J.}\ (in press).

\bibitem{Bardeen1973}
Bardeen, J.~M., Press, W.~H., and Teukolsky, S.~A. 1972, \textit{Ap.\ J.}\ (in press).

\bibitem{Berestetskii1966}
Berestetskii, V.~B. 1966, Chapter 8 of \textit{Quasistellar Sources and
Gravitational Collapse}, I.~Robinson, A.~Schild and E.~L.~Sch\"ucking eds.\ (Chicago:
University of Chicago Press).

\bibitem{Bisnovatyi1968}
Bisnovatyi-Kogan, G.~S., Zel'dovich, Ya.~B., and Sunyaev, R.~A. 1971, \textit{Soviet Astron.---A.~J.}\
\textbf{48}, 24.

\bibitem{Bondi1952}
Bondi, H. 1952, \textit{Mon.\ Not.\ Roy.\ Astron.\ Soc.}\ \textbf{112}, 195.

\bibitem{Bondi1944}
Bondi, H. and Hoyle, F. 1944, \textit{Mon.\ Not.\ Roy.\ Astron.\ Soc.}\ \textbf{104}, 273.

\bibitem{Bruenn1972}
Bruenn, S. 1972, \textit{Comments Ap.\ Space Phys.}\ \textbf{34}, 507.

\bibitem{Burbidge1972}
Burbidge, G.~R. 1972, \textit{Comments Ap.\ Space Phys.}\ (in press).

\bibitem{Chandrasekhar1939}
Chandrasekhar, S. 1939, \textit{Stellar Structure} (Chicago: University of Chicago Press);
reprinted 1957 (New York: Dover Publications, Inc.).

\bibitem{Chandrasekhar1960}
Chandrasekhar, S. 1960, \textit{Radiative Transfer} (New York: Dover Publications, Inc.).

\bibitem{Cowling1965}
Cowling, T.~G. 1965, Chapter 8 of \textit{Stellar Structure},
L.~H.~Aller and D.~B.~McLaughlin eds.\
(Chicago: University of Chicago Press).

\bibitem{Eddington1926}
Eddington, A.~S. 1926, \textit{The Internal Constitution of the Stars} (Cambridge: Cambridge
University Press).

\bibitem{Ekers1970}
Ekers, R. and Lynden-Bell, D. 1971, \textit{Astrophys.\ Lett.}\ \textbf{9}, 189.

\bibitem{Ellis1971}
Ellis, G.~F.~R. 1971, in \textit{Relativity and Cosmology},
R.~Sachs ed.\ (New York: Academic Press).

\bibitem{Felten1972}
Felten, J.~E. 1968, \textit{Ap.\ Space Sci.}\ \textbf{2}, 96.

\bibitem{Ginzburg1967}
Ginzburg, V.~L. 1967, in \textit{High Energy Astrophysics, Vol.\ 1},
C.~DeWitt, E.~Schatzman, P.~V\'eron eds.\ (New York: Gordon and Breach), p.~19.

\bibitem{Gold1965}
Gold, T., and Hoyle, F. 1965, Chapter 8 of \textit{Quasistellar Sources and
Gravitational Collapse}, I.~Robinson, A.~Schild and E.~L.~Sch\"ucking eds.\ (Chicago:
University of Chicago Press).

\bibitem{Grassberg1969}
Grassberg, E.~K. and Nadyozhin, D.~K. 1969, \textit{Ap.\ Space Sci.}\ \textbf{3}, 464.

\bibitem{Green1959}
Green, J. 1959, \textit{Ap.\ J.}\ \textbf{130}, 693.

\bibitem{Hansen1964}
Hansen, C.~J. and Matsuda, T. 1964, \textit{Canadian J.\ Phys.}

\bibitem{Hayakawa1964}
Hayakawa, S. and Matsuoka, M. 1964, \textit{Prog.\ Theor.\ Phys.\ Suppl.}\ \textbf{30}, 204.

\bibitem{Heitler1954}
Heitler, W. 1954, \textit{The Quantum Theory of Radiation} (Oxford: Clarendon Press).

\bibitem{Hoyle1939}
Hoyle, F. and Lyttleton, R.~A. 1939, \textit{Proc.\ Camb.\ Phil.\ Soc.}\ \textbf{35}, 405.

\bibitem{Hoyle1964}
Hoyle, F., Fowler, W.~A., Burbidge, G.~R. and Burbidge, E.~M. 1964, \textit{Ap.\ J.}\ \textbf{139}, 909.

\bibitem{Illarionov1972a}
Illarionov, A.~F. and Sunyaev, R.~A. 1972, \textit{Soviet Astron.---A.~J.}\ \textbf{49}, 58.

\bibitem{Illarionov1972b}
Illarionov, A.~F. and Sunyaev, R.~A. 1972, \textit{Ap.\ J.}\ (in preparation).

\bibitem{Ivanov1969}
Ivanov, L.~N., Imshennnik, V.~S. and Nadezhin, O.~K., \textit{Sci.\ Inf.\ Astr.\ Council USSR Acad.\
Sci.}\ \textbf{13}.

\bibitem{Jackson1962}
Jackson, J.~D. 1962, \textit{Classical Electrodynamics} (New York: John Wiley and Sons, Inc.).

\bibitem{Kaplan1966}
Kaplan, S.~A. and Pikel'ner, S.~B. 1970, \textit{The Interstellar Medium} (Cambridge: Harvard
University Press).

\bibitem{Kaplan1966b}
Kaplan, S.~A. 1966, \textit{Interstellar Gas Dynamics} (Oxford: Pergamon Press).

\bibitem{Kaplan1971}
Kaplan, S.~A. and Serotovich, 1971, \textit{Planetnaya Astrofizika} (Moscow: Izdatel'stvo Nauka).

\bibitem{Kompaneets1957}
Kompaneets, A.~S. 1957, \textit{Soviet Phys.\ JETP}\ \textbf{4}, 730.

\bibitem{Korchak1961}
Korchak, A.~A. and Syrovatskii, S.~I. 1961, \textit{Ap.\ J.\ Suppl.}\ \textbf{6}, 167.

\bibitem{Landau1959}
Landau, L.~D. and Lifshitz, E.~M. 1959, \textit{Fluid Mechanics} (London: Pergamon Press).

\bibitem{Leighton1959}
Leighton, R.~B. 1959, \textit{Principles of Modern Physics} (New York: McGraw-Hill).

\bibitem{Lichnerowicz1967}
Lichnerowicz, A. 1967, \textit{Relativistic Hydrodynamics and Magnetohydrodynamics} (New
York: W.~A.~Benjamin, Inc.).

\bibitem{Lichnerowicz1970}
Lichnerowicz, A. 1970, Lectures in \textit{Relativistic Fluid Dynamics}, proceedings of July 1970
(Rome: Edizioni Cremonese).

\bibitem{Lichnerowicz1971}
Lichnerowicz, A. 1971, \textit{Physica Scripta}\ \textbf{2}, 221.

\bibitem{Lindquist1966}
Lindquist, R.~W. 1966, \textit{Ann.\ Phys.}\ (N.Y.) \textbf{37}, 487.

\bibitem{LyndenBell1969}
Lynden-Bell, D. 1969, \textit{Nature}\ \textbf{223}, 690.

\bibitem{LyndenBell1971}
Lynden-Bell, D. and Rees, M.~J. 1971, \textit{Mon.\ Not.\ Roy.\ Astron.\ Soc.}\ \textbf{152}, 461.

\bibitem{Mishalas1970}
Mihalas, D. 1970, \textit{Stellar Atmospheres} (San Francisco: W.~H.~Freeman and Co.).

\bibitem{Misner1973}
Misner, C.~W., Thorne, K.~S. and Wheeler, J.~A. 1973, \textit{Gravitation} (San Francisco: W.~H.
Freeman and Co.), cited as MTW.

\bibitem{Neumann1965}
Ne'eman, Y. 1965, \textit{Ap.\ J.}\ \textbf{141}, 1303.

\bibitem{Novikov1964}
Novikov, I.~D. 1964, \textit{Astron.\ Zhur.}\ \textbf{41}, 1075.

\bibitem{Novikov1966}
Novikov, I.~D., Polnarev, A.~G. and Sunyaev, R.~A. 1973, in preparation.

\bibitem{Novikov1973a}
Novikov, I.~D. and Thorne, K.~S. 1973, in \textit{Black Holes}, C.~DeWitt and B.~S.~DeWitt eds.\
(New York: Gordon and Breach), p.~343--450.

\bibitem{Pacholczyk1970}
Pacholczyk, A.~G. 1970, \textit{Radio Astrophysics} (San Francisco: W.~H.~Freeman and Co.).

\bibitem{Paczynski1971}
Paczy\'nski, B. 1971, \textit{Ann.\ Rev.\ Astron.\ Ap.}\ \textbf{9}, 183.

\bibitem{Peebles1972}
Peebles, P.~J.~E. 1971, \textit{Gen.\ Rel.\ Grav.}\ \textbf{3}, 63.

\bibitem{Pikelner1961}
Pikel'ner, S.~B. 1961, \textit{Foundations of Cosmic Electrodynamics} (Moscow: Fizmatgiz) (in
Russian).

\bibitem{Prendergast1968}
Prendergast, K.~H. and Burbidge, G.~R. 1968, \textit{Ap.\ J.\ Letters}\ \textbf{151}, L83.

\bibitem{Pringle1972}
Pringle, J.~E. 1964, \textit{Ap.\ J.}\ \textbf{140}, 796.

\bibitem{Salpeter1964}
Salpeter, E.~N. 1964, \textit{Ap.\ J.}\ \textbf{140}, 796.

\bibitem{Shakura1972}
Shakura, N.~I. 1972, \textit{Soviet Astronomy}\ \textbf{16}.

\bibitem{Shakura1972b}
Shakura, N.~I. 1972b, \textit{Soviet Astron.---A.~J.}\ \textbf{49}, 921.

\bibitem{Shakura1973}
Shakura, N.~I. and Sunyaev, R.~A. 1973, \textit{Astron.\ Ap.}\ \textbf{24}, 337.

\bibitem{Shapiro1973a}
Shapiro, S. 1973a, \textit{Ap.\ J.}\ \textbf{180}, 531.

\bibitem{Shapiro1973b}
Shapiro, S. 1973b, \textit{Ap.\ J.}\ (in press).

\bibitem{Shvartsman1971}
Shvartsman, V.~F. 1971, \textit{Soviet Astron.---A.~J.}\ \textbf{15}, 377.

\end{thebibliography}
